% ============================================================================
% FIGURE 1: Knowledge → Activation → Output (Main Paper Figure)
% ============================================================================
% Status: FINAL - Paper-ready, theoretically correct, visually clear
% Date: 2025-12-28
% Usage: % ============================================================================
% FIGURE 1: Knowledge → Activation → Output (Main Paper Figure)
% ============================================================================
% Status: FINAL - Paper-ready, theoretically correct, visually clear
% Date: 2025-12-28
% Usage: % ============================================================================
% FIGURE 1: Knowledge → Activation → Output (Main Paper Figure)
% ============================================================================
% Status: FINAL - Paper-ready, theoretically correct, visually clear
% Date: 2025-12-28
% Usage: % ============================================================================
% FIGURE 1: Knowledge → Activation → Output (Main Paper Figure)
% ============================================================================
% Status: FINAL - Paper-ready, theoretically correct, visually clear
% Date: 2025-12-28
% Usage: \input{ESL_Figure_1_Main.tex} in main paper
% ============================================================================

% Required in preamble:
% \usepackage{tikz}
% \usetikzlibrary{arrows.meta,positioning,shapes.geometric,calc}

% ============================================================================
% FIGURE 1: Main Diagram
% ============================================================================

\begin{figure}[t]
\centering
\begin{tikzpicture}[
    node distance=2.0cm,
    block/.style={
        rectangle,
        draw=black,
        thick,
        minimum width=3.5cm,
        minimum height=1.6cm,
        align=center,
        rounded corners=4pt
    },
    prompt/.style={
        rectangle,
        draw=black,
        thick,
        dashed,
        minimum width=2.5cm,
        minimum height=1.2cm,
        align=center,
        rounded corners=4pt,
        fill=gray!10
    },
    arrow/.style={-{Stealth[length=3.5mm]}, thick},
    label/.style={font=\small\itshape, text=black!70}
]

% =========================================
% Block A: Latent Knowledge (top)
% =========================================
\node[block, fill=blue!12] (M) {
    \textbf{Latent Knowledge}\\[3pt]
    $M_{\text{latent}}(x)$
};
\node[label, right=0.3cm of M, align=left] {
    available\\
    but inactive
};

% =========================================
% Block B: Activation (middle) - THE BOTTLENECK
% =========================================
\node[block, fill=green!15, below=of M] (alpha) {
    \textbf{Knowledge Activation}\\[3pt]
    $\alpha := \mathbb{P}(A(x) \mid M_{\text{latent}}=1)$
};

% Bottleneck annotation
\node[right=0.3cm of alpha, align=left, font=\small] (bottleneck) {
    \textcolor{red!70!black}{\textbf{epistemic}}\\
    \textcolor{red!70!black}{\textbf{bottleneck}}
};

% =========================================
% Block C: Generative Output (bottom)
% =========================================
\node[block, fill=orange!12, below=of alpha] (output) {
    \textbf{Generative Output}\\[3pt]
    Decision / Text
};
\node[label, right=0.3cm of output, align=left] {
    applied /\\
    not applied
};

% =========================================
% Prompt node (acts on activation, NOT knowledge)
% =========================================
\node[prompt, left=2.5cm of alpha] (pi) {
    \textbf{Prompt}\\[2pt]
    $\pi$
};

% =========================================
% Arrows
% =========================================
% Main vertical flow
\draw[arrow] (M) -- (alpha);
\draw[arrow] (alpha) -- (output);

% Prompt to Activation (dashed - modulates, not creates)
\draw[arrow, dashed, gray!70] (pi) -- (alpha)
    node[midway, above, font=\footnotesize\itshape] {modulates};

% =========================================
% Key insight box
% =========================================
\node[draw=red!60!black, thick, rounded corners=3pt, fill=red!5,
      below=0.8cm of output, align=center, font=\small] (insight) {
    \textbf{Key Finding:} $M_{\text{latent}}(x) = 1 \;\not\Rightarrow\; A(x)$
};

\end{tikzpicture}

\caption{\textbf{Functional decomposition of language model behavior.} Latent knowledge (availability) does not directly determine generative output. Instead, knowledge must pass through an activation stage, characterized by activation competence $\alpha$, which governs whether available knowledge is applied during generation. Prompting and epistemic constraints act on this activation stage but do not modify latent knowledge itself. This is a functional abstraction, not an architectural claim.}
\label{fig:activation-model}
\end{figure}


% ============================================================================
% COMPONENT DEFINITIONS (for reference in text)
% ============================================================================
%
% (A) Latent Knowledge — Availability
%     M_latent(x) = 1 iff model can answer direct query about x
%     - not visible in generative output
%     - not equivalent to correct behavior
%     - no guarantee of application
%
% (B) Knowledge Activation — the bottleneck
%     α := P(A(x) | M_latent(x) = 1)
%     - independent competence
%     - bounded prompt-controllability
%     - model-dependent
%     - epistemically critical
%
% (C) Generative Output — Behavior
%     Output ~ f_θ(x | α)
%     - classification, text, decision, assertion
%     - fluency can dominate when α is low
%
% ============================================================================
% WHAT THIS FIGURE DOES NOT CLAIM (important for reviewers)
% ============================================================================
%
% ❌ no neural module
% ❌ no architectural assumption
% ❌ no sequential processing claim
% ✓  purely functional decomposition
%
% The diagram is an abstraction, not an implementation.
%
 in main paper
% ============================================================================

% Required in preamble:
% \usepackage{tikz}
% \usetikzlibrary{arrows.meta,positioning,shapes.geometric,calc}

% ============================================================================
% FIGURE 1: Main Diagram
% ============================================================================

\begin{figure}[t]
\centering
\begin{tikzpicture}[
    node distance=2.0cm,
    block/.style={
        rectangle,
        draw=black,
        thick,
        minimum width=3.5cm,
        minimum height=1.6cm,
        align=center,
        rounded corners=4pt
    },
    prompt/.style={
        rectangle,
        draw=black,
        thick,
        dashed,
        minimum width=2.5cm,
        minimum height=1.2cm,
        align=center,
        rounded corners=4pt,
        fill=gray!10
    },
    arrow/.style={-{Stealth[length=3.5mm]}, thick},
    label/.style={font=\small\itshape, text=black!70}
]

% =========================================
% Block A: Latent Knowledge (top)
% =========================================
\node[block, fill=blue!12] (M) {
    \textbf{Latent Knowledge}\\[3pt]
    $M_{\text{latent}}(x)$
};
\node[label, right=0.3cm of M, align=left] {
    available\\
    but inactive
};

% =========================================
% Block B: Activation (middle) - THE BOTTLENECK
% =========================================
\node[block, fill=green!15, below=of M] (alpha) {
    \textbf{Knowledge Activation}\\[3pt]
    $\alpha := \mathbb{P}(A(x) \mid M_{\text{latent}}=1)$
};

% Bottleneck annotation
\node[right=0.3cm of alpha, align=left, font=\small] (bottleneck) {
    \textcolor{red!70!black}{\textbf{epistemic}}\\
    \textcolor{red!70!black}{\textbf{bottleneck}}
};

% =========================================
% Block C: Generative Output (bottom)
% =========================================
\node[block, fill=orange!12, below=of alpha] (output) {
    \textbf{Generative Output}\\[3pt]
    Decision / Text
};
\node[label, right=0.3cm of output, align=left] {
    applied /\\
    not applied
};

% =========================================
% Prompt node (acts on activation, NOT knowledge)
% =========================================
\node[prompt, left=2.5cm of alpha] (pi) {
    \textbf{Prompt}\\[2pt]
    $\pi$
};

% =========================================
% Arrows
% =========================================
% Main vertical flow
\draw[arrow] (M) -- (alpha);
\draw[arrow] (alpha) -- (output);

% Prompt to Activation (dashed - modulates, not creates)
\draw[arrow, dashed, gray!70] (pi) -- (alpha)
    node[midway, above, font=\footnotesize\itshape] {modulates};

% =========================================
% Key insight box
% =========================================
\node[draw=red!60!black, thick, rounded corners=3pt, fill=red!5,
      below=0.8cm of output, align=center, font=\small] (insight) {
    \textbf{Key Finding:} $M_{\text{latent}}(x) = 1 \;\not\Rightarrow\; A(x)$
};

\end{tikzpicture}

\caption{\textbf{Functional decomposition of language model behavior.} Latent knowledge (availability) does not directly determine generative output. Instead, knowledge must pass through an activation stage, characterized by activation competence $\alpha$, which governs whether available knowledge is applied during generation. Prompting and epistemic constraints act on this activation stage but do not modify latent knowledge itself. This is a functional abstraction, not an architectural claim.}
\label{fig:activation-model}
\end{figure}


% ============================================================================
% COMPONENT DEFINITIONS (for reference in text)
% ============================================================================
%
% (A) Latent Knowledge — Availability
%     M_latent(x) = 1 iff model can answer direct query about x
%     - not visible in generative output
%     - not equivalent to correct behavior
%     - no guarantee of application
%
% (B) Knowledge Activation — the bottleneck
%     α := P(A(x) | M_latent(x) = 1)
%     - independent competence
%     - bounded prompt-controllability
%     - model-dependent
%     - epistemically critical
%
% (C) Generative Output — Behavior
%     Output ~ f_θ(x | α)
%     - classification, text, decision, assertion
%     - fluency can dominate when α is low
%
% ============================================================================
% WHAT THIS FIGURE DOES NOT CLAIM (important for reviewers)
% ============================================================================
%
% ❌ no neural module
% ❌ no architectural assumption
% ❌ no sequential processing claim
% ✓  purely functional decomposition
%
% The diagram is an abstraction, not an implementation.
%
 in main paper
% ============================================================================

% Required in preamble:
% \usepackage{tikz}
% \usetikzlibrary{arrows.meta,positioning,shapes.geometric,calc}

% ============================================================================
% FIGURE 1: Main Diagram
% ============================================================================

\begin{figure}[t]
\centering
\begin{tikzpicture}[
    node distance=2.0cm,
    block/.style={
        rectangle,
        draw=black,
        thick,
        minimum width=3.5cm,
        minimum height=1.6cm,
        align=center,
        rounded corners=4pt
    },
    prompt/.style={
        rectangle,
        draw=black,
        thick,
        dashed,
        minimum width=2.5cm,
        minimum height=1.2cm,
        align=center,
        rounded corners=4pt,
        fill=gray!10
    },
    arrow/.style={-{Stealth[length=3.5mm]}, thick},
    label/.style={font=\small\itshape, text=black!70}
]

% =========================================
% Block A: Latent Knowledge (top)
% =========================================
\node[block, fill=blue!12] (M) {
    \textbf{Latent Knowledge}\\[3pt]
    $M_{\text{latent}}(x)$
};
\node[label, right=0.3cm of M, align=left] {
    available\\
    but inactive
};

% =========================================
% Block B: Activation (middle) - THE BOTTLENECK
% =========================================
\node[block, fill=green!15, below=of M] (alpha) {
    \textbf{Knowledge Activation}\\[3pt]
    $\alpha := \mathbb{P}(A(x) \mid M_{\text{latent}}=1)$
};

% Bottleneck annotation
\node[right=0.3cm of alpha, align=left, font=\small] (bottleneck) {
    \textcolor{red!70!black}{\textbf{epistemic}}\\
    \textcolor{red!70!black}{\textbf{bottleneck}}
};

% =========================================
% Block C: Generative Output (bottom)
% =========================================
\node[block, fill=orange!12, below=of alpha] (output) {
    \textbf{Generative Output}\\[3pt]
    Decision / Text
};
\node[label, right=0.3cm of output, align=left] {
    applied /\\
    not applied
};

% =========================================
% Prompt node (acts on activation, NOT knowledge)
% =========================================
\node[prompt, left=2.5cm of alpha] (pi) {
    \textbf{Prompt}\\[2pt]
    $\pi$
};

% =========================================
% Arrows
% =========================================
% Main vertical flow
\draw[arrow] (M) -- (alpha);
\draw[arrow] (alpha) -- (output);

% Prompt to Activation (dashed - modulates, not creates)
\draw[arrow, dashed, gray!70] (pi) -- (alpha)
    node[midway, above, font=\footnotesize\itshape] {modulates};

% =========================================
% Key insight box
% =========================================
\node[draw=red!60!black, thick, rounded corners=3pt, fill=red!5,
      below=0.8cm of output, align=center, font=\small] (insight) {
    \textbf{Key Finding:} $M_{\text{latent}}(x) = 1 \;\not\Rightarrow\; A(x)$
};

\end{tikzpicture}

\caption{\textbf{Functional decomposition of language model behavior.} Latent knowledge (availability) does not directly determine generative output. Instead, knowledge must pass through an activation stage, characterized by activation competence $\alpha$, which governs whether available knowledge is applied during generation. Prompting and epistemic constraints act on this activation stage but do not modify latent knowledge itself. This is a functional abstraction, not an architectural claim.}
\label{fig:activation-model}
\end{figure}


% ============================================================================
% COMPONENT DEFINITIONS (for reference in text)
% ============================================================================
%
% (A) Latent Knowledge — Availability
%     M_latent(x) = 1 iff model can answer direct query about x
%     - not visible in generative output
%     - not equivalent to correct behavior
%     - no guarantee of application
%
% (B) Knowledge Activation — the bottleneck
%     α := P(A(x) | M_latent(x) = 1)
%     - independent competence
%     - bounded prompt-controllability
%     - model-dependent
%     - epistemically critical
%
% (C) Generative Output — Behavior
%     Output ~ f_θ(x | α)
%     - classification, text, decision, assertion
%     - fluency can dominate when α is low
%
% ============================================================================
% WHAT THIS FIGURE DOES NOT CLAIM (important for reviewers)
% ============================================================================
%
% ❌ no neural module
% ❌ no architectural assumption
% ❌ no sequential processing claim
% ✓  purely functional decomposition
%
% The diagram is an abstraction, not an implementation.
%
 in main paper
% ============================================================================

% Required in preamble:
% \usepackage{tikz}
% \usetikzlibrary{arrows.meta,positioning,shapes.geometric,calc}

% ============================================================================
% FIGURE 1: Main Diagram
% ============================================================================

\begin{figure}[t]
\centering
\begin{tikzpicture}[
    node distance=2.0cm,
    block/.style={
        rectangle,
        draw=black,
        thick,
        minimum width=3.5cm,
        minimum height=1.6cm,
        align=center,
        rounded corners=4pt
    },
    prompt/.style={
        rectangle,
        draw=black,
        thick,
        dashed,
        minimum width=2.5cm,
        minimum height=1.2cm,
        align=center,
        rounded corners=4pt,
        fill=gray!10
    },
    arrow/.style={-{Stealth[length=3.5mm]}, thick},
    label/.style={font=\small\itshape, text=black!70}
]

% =========================================
% Block A: Latent Knowledge (top)
% =========================================
\node[block, fill=blue!12] (M) {
    \textbf{Latent Knowledge}\\[3pt]
    $M_{\text{latent}}(x)$
};
\node[label, right=0.3cm of M, align=left] {
    available\\
    but inactive
};

% =========================================
% Block B: Activation (middle) - THE BOTTLENECK
% =========================================
\node[block, fill=green!15, below=of M] (alpha) {
    \textbf{Knowledge Activation}\\[3pt]
    $\alpha := \mathbb{P}(A(x) \mid M_{\text{latent}}=1)$
};

% Bottleneck annotation
\node[right=0.3cm of alpha, align=left, font=\small] (bottleneck) {
    \textcolor{red!70!black}{\textbf{epistemic}}\\
    \textcolor{red!70!black}{\textbf{bottleneck}}
};

% =========================================
% Block C: Generative Output (bottom)
% =========================================
\node[block, fill=orange!12, below=of alpha] (output) {
    \textbf{Generative Output}\\[3pt]
    Decision / Text
};
\node[label, right=0.3cm of output, align=left] {
    applied /\\
    not applied
};

% =========================================
% Prompt node (acts on activation, NOT knowledge)
% =========================================
\node[prompt, left=2.5cm of alpha] (pi) {
    \textbf{Prompt}\\[2pt]
    $\pi$
};

% =========================================
% Arrows
% =========================================
% Main vertical flow
\draw[arrow] (M) -- (alpha);
\draw[arrow] (alpha) -- (output);

% Prompt to Activation (dashed - modulates, not creates)
\draw[arrow, dashed, gray!70] (pi) -- (alpha)
    node[midway, above, font=\footnotesize\itshape] {modulates};

% =========================================
% Key insight box
% =========================================
\node[draw=red!60!black, thick, rounded corners=3pt, fill=red!5,
      below=0.8cm of output, align=center, font=\small] (insight) {
    \textbf{Key Finding:} $M_{\text{latent}}(x) = 1 \;\not\Rightarrow\; A(x)$
};

\end{tikzpicture}

\caption{\textbf{Functional decomposition of language model behavior.} Latent knowledge (availability) does not directly determine generative output. Instead, knowledge must pass through an activation stage, characterized by activation competence $\alpha$, which governs whether available knowledge is applied during generation. Prompting and epistemic constraints act on this activation stage but do not modify latent knowledge itself. This is a functional abstraction, not an architectural claim.}
\label{fig:activation-model}
\end{figure}


% ============================================================================
% COMPONENT DEFINITIONS (for reference in text)
% ============================================================================
%
% (A) Latent Knowledge — Availability
%     M_latent(x) = 1 iff model can answer direct query about x
%     - not visible in generative output
%     - not equivalent to correct behavior
%     - no guarantee of application
%
% (B) Knowledge Activation — the bottleneck
%     α := P(A(x) | M_latent(x) = 1)
%     - independent competence
%     - bounded prompt-controllability
%     - model-dependent
%     - epistemically critical
%
% (C) Generative Output — Behavior
%     Output ~ f_θ(x | α)
%     - classification, text, decision, assertion
%     - fluency can dominate when α is low
%
% ============================================================================
% WHAT THIS FIGURE DOES NOT CLAIM (important for reviewers)
% ============================================================================
%
% ❌ no neural module
% ❌ no architectural assumption
% ❌ no sequential processing claim
% ✓  purely functional decomposition
%
% The diagram is an abstraction, not an implementation.
%
